\section{Confound analysis in full}
\label{app:confounds}

Table~\ref{tab:confounds} summarises; this appendix gives the operational detail
for the controls that are not self-explanatory.

\paragraph{Mechanical pooling shift.} Three independent answers, in order of
strength. (i) Where per-cell consensus is full and the per-cell margin is
positive everywhere, Proposition~\ref{prop:deletion} excludes the confound
formally. We report consensus, agreement with the emitted symbol, and the minimum
per-cell margin before and after damage at every $\Delta$, so a reader can see
where the proposition applies. (ii) The frozen-state deletion control of
Section~\ref{sec:mechanical}. (iii) Evaluating the same states under sum-pool,
median-pool, and per-cell majority vote: if the phenomenon exists only under mean
pooling it is an artifact.

The pooling choice is not neutral, and the reason is worth stating. Under full
consensus, mean-pool argmax is deletion-invariant but the logit \emph{magnitude}
barely moves. Under sum-pool with dead cells contributing zero, deleting cells
shrinks all logits, flattening the softmax and raising $\Ent(M|R)$ even when the
argmax is stable. The pooling choice therefore determines whether damage
manifests as a change of \emph{meaning} or a change of \emph{confidence} --- so
we train with mean-pool and report the sum-pool readout on the same states as a
diagnostic that separates the two.

\paragraph{Confidence versus identity.} A sender may become less certain without
changing its preferred symbol, which lowers sampled accuracy through a flatter
softmax rather than through a changed code. We report argmax accuracy and sampled
accuracy at every $\Delta$, decomposing the total change into a drift component
(argmax flips) and an uncertainty component (the gap between the two), alongside
the logit margin and $\Ent(M|R)$. A temperature sweep settles it: if the drop
vanishes as $\tau \to 0$, it was confidence.

\paragraph{Sensor destruction.} If the damage disk overlaps the sensor patch the
referent information is literally deleted, which is trivial rather than
interesting. We stratify every episode by sensor overlap and headline the
non-overlapping stratum, and run sensor-targeted damage separately as an explicit
contrast between destroying the input pathway and destroying the computation.

\paragraph{Mid-regeneration.} ``The sender simply was not finished'' is excluded
two ways: by running to $\Delta = 256 \gg T_{\mathrm{grow}} = 64$ and reporting
the asymptote, and by reporting semantic metrics binned by IoU rather than by
wall-clock, which definitionally controls for how far regeneration has
progressed. We note the hazard in the second: binning on a post-treatment
variable is a collider, so severity distributions are reported alongside, and the
matched analysis is treated as supporting rather than primary.

\paragraph{Training-condition confounds.} Arms that differ in training may differ
in pre-damage accuracy, contaminating every comparison. Three layers: ANCOVA with
pre-damage accuracy as a covariate; checkpoint matching to within 2 points; and a
third arm trained with damage in the pool but with the task loss detached during
damaged steps, which learns to regrow but never learns to communicate while
damaged. That third arm is what makes the robustness claim mechanistic rather
than correlational.

\section{Sampling and cost}
\label{app:sampling}

Referents are enumerated rather than sampled. Damage placements are the dominant
variance source and the bootstrap unit: $J = 32$ placements per condition, with
the \emph{same} placement seeds reused across training conditions and damage
magnitudes, since pairing is free and is the difference between a significant and
a non-significant comparison at 8 seeds. Episodes per condition per seed are
$N \times J \times K = 8 \times 32 \times 4 = 1024$.

Evaluation is dense early and logarithmic later,
$\Delta \in \{0,1,2,4,8,12,16,24,32,48,64,96,128,192,256\}$. At each $\Delta$ the
receiver is run fresh from its own seed using the symbol emitted at that
$\Delta$, which decouples receiver dynamics from sender time; the alternative --- a
single continuous receiver rollout with a time-varying symbol --- entangles them
and is materially harder to interpret.

\section{Hardware and cost}
\label{app:hardware}

All experiments run on a single Apple M2 with 8\,GB of unified memory. Gradient
checkpointing every 8 steps is required rather than optional: activation memory
at $T=64$, batch 32 falls from 3300\,MB to 1110\,MB, and is then nearly flat in
rollout length (1119\,MB at $T=128$). That flatness is what makes long rollouts
affordable, which matters because information advances roughly half a cell per
step under stochastic updates, so a short rollout would leave most of the body
blind to the referent.

The morphology run reported in Section~\ref{sec:results-morph} took 94 minutes
for 6000 steps.
